\begin{textsample}{POS Dim 2 – rr – Score 25.00 – rr/rr\_em\_57.txt}  \label{ex:f2_pos_150}
5.9.4 A \textbf{oferta} de formação \textbf{técnica} e profissional está \textbf{prevista} na meta 11 do Plano Nacional de Educação – PNE ( BRASIL, 2014 ) e Plano Estadual de Educação de Roraima – PEE/RR ( RORAIMA, 2015 ), aprovados para o decênio ( 2014-2024 ), que \textbf{prevê} triplicar as matrículas na EPT de nível médio, assegurando a qualidade da \textbf{oferta} e pelo menos 50 \% da expansão no segmento público.
Esta meta, está alinhada ao desejo dos jovens roraimenses, constatada pela escuta online, quando 27 \% dos jovens responderam que 1 ( uma ) hora a mais por dia na escola, gostariam de receber atividades de preparação profissional.
O Art. 7º da Resolução CNE nº 6 ( BRASIL, 2012c ), a \textbf{oferta} de EPT pelo sistema estadual se dão nas formas articulada e subsequente ao Ensino Médio, conforme quadro abaixo : | MODALIDADE | CATEGORIAS | CONDIÇÃO | |---|---|---| | Articulada | Integrada | \textbf{Ofertada} somente a quem já tenha concluído o Ensino Fundamental, com matrícula única na mesma \textbf{instituição}, de modo a conduzir o estudante à \textbf{habilitação} profissional \textbf{técnica} de nível médio ao mesmo tempo em que conclui a última etapa da Educação Básica. | | Articulada | Concomitante | \textbf{Ofertada} a quem ingressa no Ensino Médio ou já o estejam \textbf{cursando}, \textbf{efetuando} -se matrículas distintas para cada \textbf{curso}, aproveitando oportunidades educacionais disponíveis, seja em unidades de ensino da mesma \textbf{instituição} ou em distintas \textbf{instituições} de ensino. | | Articulada | Concomitante na forma | \textbf{Oferta} desenvolvida simultaneamente em distintas \textbf{instituições} educacionais, mas integrada no conteúdo, mediante a ação de convênio ou acordo de intercomplementaridade, para a execução de projeto pedagógico unificado. | | Subsequente | | Desenvolvida em \textbf{cursos} \textbf{destinados} exclusivamente a quem já tenha concluído o Ensino Médio. | Em consonância com a Lei 13.415 ( BRASIL, 2017a ), o sistema de ensino estadual, sintetiza a organização da \textbf{oferta} de EPT, conforme se segue : - Aluno concluinte pode \textbf{cursar} mais de um \textbf{itinerário} ; - Possibilidade de ser organizado em módulos 32 e adotar o sistema de créditos com terminalidade específica ; 32 O sistema modular é uma modalidade de ensino que garante o EPT em \textbf{localidades} mais distantes, que precisavam de atendimento imediato, porém não é possível construir uma escola com toda a estrutura do ensino regular por ter menos alunos. - As escolas deverão orientar os alunos no processo de escolha das áreas de conhecimento ou de atuação profissional ; - As \textbf{instituições} de ensino emitirão certificado com validade nacional ; - Possibilidade de reconhecer competências e firmar convênios com \textbf{instituições} de educação a distância com \textbf{notório} reconhecimento ; Resumidamente, em acordo às \textbf{Diretrizes} da Resolução CNE nº 3 ( BRASIL, 2018d ), o sistema de ensino estadual, considera a \textbf{oferta} do \textbf{itinerário} formativo EPT, como : - a inclusão de vivências práticas de trabalho, constante de \textbf{carga} \textbf{horária} específica, no setor produtivo ou em ambientes de simulação, estabelecendo parcerias e fazendo uso, quando aplicável, de instrumentos estabelecidos pela \textbf{legislação} sobre aprendizagem profissional ; - a possibilidade de concessão de certificados intermediários de \textbf{qualificação} para o trabalho, quando a formação for estruturada e organizada em etapas com terminalidade específica. - os profissionais com \textbf{notório} saber, são reconhecidos pelo sistema de ensino local, podendo atuar como docentes no Ensino Médio, apenas no \textbf{itinerário} EPT para ministrar conteúdos afins à sua formação ou experiência profissional, devidamente comprovadas, conforme inciso IV do \textbf{art}. 61 da LDB.
A docência nas \textbf{instituições} que \textbf{ofertam} o \textbf{itinerário} EPT será realizada por profissionais com comprovada competência \textbf{técnica} referente ao saber operativo de atividades inerentes à respectiva formação \textbf{técnica} e profissional.
O sistema de ensino local opta por estabelecer parcerias com outras \textbf{instituições} de ensino para \textbf{oferta} de diferentes \textbf{itinerários}.
No caso da formação \textbf{técnica} e profissional, mesmo estudantes que não optarem inicialmente por esse \textbf{itinerário} podem realizar \textbf{cursos} \textbf{técnicos} ou FICs em escolas credenciadas da região.
O CEE \textbf{normatiza} os critérios para o estabelecimento de parcerias, sendo a \textbf{instituição} de origem do aluno a responsável por estabelecer \textbf{diretrizes} para o acompanhamento dos \textbf{cursos} realizados em outras \textbf{instituições}.
O roteiro de organização de trajetórias para a Produção Curricular específica do \textbf{itinerário} formativo EPT, considera : 1.
Tipo de \textbf{oferta} de Itinerário Formativo de Formação \textbf{Técnica} Profissional \textbf{Curso} \textbf{Técnico} ; FICs ; Programa de Aprendizagem ( Projeto Aprendiz ). 2.
Matrizes de \textbf{Cursos} \textbf{Técnicos} e Planos de \textbf{Curso} : As Matrizes curriculares de \textbf{Cursos} \textbf{Técnicos} serão revisadas a partir da parte comum do \textbf{currículo}, organizada por competência e incorporando componentes curriculares/unidades que contemplem os eixos estruturantes considerando foco pedagógico, competências e habilidades ). 3.
Planos de \textbf{Curso} : Conforme Art. 20 da Resolução CNE nº 6 ( BRASIL, 2012c ), os Plano de \textbf{Curso} devem apresentar : - a identificação do \textbf{curso} justificativa e objetivos ; - requisitos e formas de acesso ; - perfil profissional das saídas intermediárias e perfil profissional de conclusão ; - organização curricular ; - critérios de aproveitamento de conhecimentos e experiências anteriores ; - critérios e procedimentos de avaliação de aprendizagem ; - biblioteca, laboratórios, instalações e equipamentos ; - perfil de professores, instrutores e \textbf{técnicos} ; - certificados e diplomas a serem emitidos. 4.
Planos de \textbf{Curso} de FICs : Segundo CBO ( BRASIL, 2010 ) os Planos de \textbf{Cursos} de FICs devem conter os mesmos requisitos acima. 5.
Ementas de \textbf{eletivas} : Sugere -se que as ementas de \textbf{eletivas} contenham : tema, competências, habilidades, objetos de conhecimento, \textbf{carga} \textbf{horária}, perfil docente, quantidade de estudantes, recursos, avaliação. 6.
Ementas de Módulo de Preparação Básica para o Trabalho : As ementas do Módulo de Preparação Básica para o Trabalho, devem conter as 4 unidades dos eixos estruturantes, de 60 a 80h cada um, sendo eles : Investigação Científica ; Processos Criativos ; Mediação e Intervenção Sociocultural e Empreendedorismo, composto de : tema, competências, habilidades, objetos de conhecimento, \textbf{carga} \textbf{horária}, perfil docente, quantidade de estudantes, recursos, avaliação. 7.
Organização Curricular para participação no Projeto Aprendiz Esta organização deve apresentar uma composição de trajetórias possíveis articulando com as demais \textbf{ofertas} de \textbf{itinerários} formativos de áreas de conhecimento.

% matched lemmas: art, carga, currículo, cursar, curso, destinar, diretriz, efetuar, eletivo, habilitação, horário, instituição, itinerário, legislação, localidade, normatizar, notório, oferta, ofertar, prever, qualificação, técnico
\end{textsample}
